\section{An exact-marginal Diophantine law}\label{sec:farey}
This section starts from the $K_5$ graph-correlation nonclosure witness of Dykema, Paulsen and Prakash~\cite{DPP} and the scalar-sum and Coxeter framework of Kruglyak, Rabanovich and Samoilenko~\cite{KRS}; see also the related nonclosure construction of Musat--R{\o}rdam~\cite{MR}. We use the quantitative exact-marginal carrier bounds of~\cite[Lemmas 4.4 and 4.6]{KFive}, version 1.0.2. The new conclusions here are the uniform Farey comparison and its Diophantine endpoint curve. Its carrier is a finite direct sum of matrix algebras with a faithful tracial state. Cost is the largest matrix-block dimension. Fix an irrational
\[
 \lambda\in\left(\frac{5-\sqrt5}{2},\frac{5+\sqrt5}{2}\right),
\]
and five projections with $\tau(P_i)=\lambda/5$. Write $S=\sum_iP_i$. The deficit is $\operatorname{Var}_\tau(S)$, by the identity $\sum_{i\ne j}\tau(P_iP_j)=\tau(S^2)-\tau(S)$. Let $D_\lambda(\eps)$ be the least block cost at deficit at most $\eps$.

Let $b_N<\lambda<a_N$ be the nearest rationals with denominator at most $N$, put
\[
 g(N)=(a_N-\lambda)(\lambda-b_N),\qquad
 H(\eps)=\min\{N:g(N)\le\eps\},
\]
and let $\mu=\mu(\lambda)\in[2,\infty]$ be the irrationality exponent.

\begin{theorem}[Farey law and endpoint curve]\label{thm:farey}
For every $\eps>0$,
\[
 H(\eps)\le D_\lambda(\eps)\le10H(\eps).
\]
Its lower and upper dimension exponents are respectively
\begin{equation}\label{eq:farey-endpoints}
 \liminf_{\eps\downarrow0}\frac{\log D_\lambda(\eps)}{\log(1/\eps)}
 =\frac1{\mu+2},\qquad
 \limsup_{\eps\downarrow0}\frac{\log D_\lambda(\eps)}{\log(1/\eps)}
 =\frac{\mu-1}{2\mu},
\end{equation}
with endpoint pair $(0,1/2)$ when $\mu=\infty$. Uniformly over the parameter, $D_\lambda(\eps)\le10\lceil1/(2\sqrt\eps)\rceil$.
\end{theorem}
\begin{proof}
Suppose the carrier has block dimensions at most $N$. On block $l$, the mean $s_l=\tr_{d_l}(S_l)$ is rational with reduced denominator at most $N$. Total variance gives
\[
 \operatorname{Var}_\tau(S)\ge\sum_lw_l(s_l-\lambda)^2.
\]
No $s_l$ lies in $(b_N,a_N)$ and $\sum_lw_ls_l=\lambda$. Taking expectations of $(s_l-a_N)(s_l-b_N)\ge0$ shows that the right side is at least $g(N)$. Nonscalar blocks and arbitrarily many central summands therefore cannot improve the lower bound.

For the upper bound, Lemmas~4.4 and 4.6 of~\cite{KFive} realize any rational $p/q$ in the closed window as a scalar sum of five projections in dimension at most $2q$, and cyclic symmetrization makes ranks equal at a factor-five cost. We check coarse Farey levels as well. Put $\alpha=(5-\sqrt5)/2=[1;2,1,1,\ldots]$. Its best lower approximants under a denominator bound are $1,4/3,11/8,29/21,\ldots$, generated by
\[
 r_0=1,\qquad r_{j+1}=1+\frac1{4-r_j}.
\]
These chain points lie outside the closed scalar-sum window, so Lemma 4.6 of~\cite{KFive} does not supply them. The first is realized in dimension one. The Coxeter transformation of~\cite{KRS} sends $p/q$ in dimension $q$ to $(5q-p)/(4q-p)$ in dimension $4q-p$. Its determinant is one, so there is no cancellation. Thus every exterior lower approximant is realized in its denominator dimension. Complementation supplies the upper exterior chain. If a Farey neighbor of the interior $\lambda$ lies below $\alpha$, it is necessarily the best lower approximant to $\alpha$ at that denominator bound; a better one would contradict its being the neighbor of $\lambda$. The upper case is symmetric. Consequently every needed neighbor, including coarse ones, has the same realization bound.

At $N=H(\eps)$ realize both neighbors, symmetrize each, and put weight $(\lambda-b_N)/(a_N-b_N)$ on the upper block. The marginals are exactly $\lambda/5$ and the variance is exactly $g(N)$. Maximum block dimension is at most $10N$. Total carrier dimension is at most $20N$, so its exponents agree, although the literal factor ten belongs to the block convention.

We prove the endpoints, including irregular continued fractions. Write $p_j/q_j$ for the convergents and
\[
 M=\limsup_j\frac{\log q_{j+1}}{\log q_j}=\mu-1.
\]
The identity follows from the convergent error bounds and best approximation. For $q=q_j$, $r=q_{j-1}/q$, and complete quotient $t=[a_{j+1};a_{j+2},\ldots]$, the convergent and the opposite intermediate convergent with index $0\le h\le a_{j+1}$ have error product
\begin{equation}\label{eq:intermediate}
 g_h=\frac{t-h}{q^4(h+r)(t+r)^2}.
\end{equation}
For $h<a_{j+1}$ they are the Farey neighbors from level $\max\{q,q_{j-1}+hq\}$ until the next mediant denominator becomes available. This follows by substituting $\lambda=(p_jt+p_{j-1})/(q_jt+q_{j-1})$ and the determinant identity for consecutive convergents.

For finite $\mu$, fix $\nu>\mu$. Every reduced fraction obeys $|\lambda-p/q|\ge c_\nu q^{-\nu}$. Farey adjacency gives gap $1/(q_aq_b)\ge N^{-2}$, so one error is at least $1/(2N^2)$ and the other at least $c_\nu N^{-\nu}$. Thus $g(N)\ge(c_\nu/2)N^{-(\nu+2)}$, implying lower endpoint at least $1/(\mu+2)$. Conversely at $N=q_j+q_{j-1}$ the convergent and first opposite intermediate convergent give
\[
 g(N)\le C/(q_j^3q_{j+1}),\qquad N\le2q_j.
\]
The consecutive pair gives the same bound if the next quotient is one. Select a subsequence on which $\log q_{j+1}/\log q_j\to M$ and take the actual deficit $g(N)$. Since $H(g(N))\le N$, the lower endpoint is at most $1/(3+M)$. The same argument gives zero when $M=\infty$.

For the upper endpoint choose the largest convergent denominator $q\le N/2$ and let $Q$ be the next. Then $Q>N/2$ and $Q\le Cq^{\nu-1}$, hence $q\ge cN^{1/(\nu-1)}$. If $Q\le N$, consecutive convergents give product at most $1/(q^2Q^2)$. If $Q>N$, the last opposite intermediate denominator below $N$ is at least $N/2$, giving product at most $C/(q^2N^2)$. Therefore
\[
 g(N)\le C'N^{-2-2/(\nu-1)}.
\]
Inversion and $\nu\downarrow\mu$ bound the upper endpoint by $(\mu-1)/(2\mu)$. For $M>1$, use pairs $q,Q$ approaching $M$ and $N=\lfloor Q/2\rfloor$. Their partial quotient tends to infinity, and~\eqref{eq:intermediate} gives $g(N)\asymp(q^2Q^2)^{-1}$. At $\eps=g(N)/2$, $H(\eps)>N$, giving upper endpoint at least $M/(2(M+1))$. When $M=1$, the lower endpoint and the upper bound already agree at $1/4$.

Finally $q_a+q_b>N$ implies $q_aq_b\ge N$. The product of two nonnegative errors with sum $1/(q_aq_b)$ is at most $1/(4q_a^2q_b^2)\le1/(4N^2)$. This proves the universal bound and the missing upper bound $1/2$ when $M=\infty$.
\end{proof}

Badly approximable parameters have bounded partial quotients and $g(N)\asymp N^{-4}$, hence $D_\lambda(\eps)=\Theta(\eps^{-1/4})$. Roth's theorem~\cite{Roth} gives $\mu=2$ for algebraic irrationals, yielding $\eps^{-1/4+o(1)}$, without asserting a uniform $\Theta$ law for every algebraic parameter. The standard continued-fraction construction~\cite[Chapter 1]{Bugeaud}, choosing partial quotients with prescribed $\limsup\log q_{j+1}/\log q_j$ realizes every $\mu\in[2,\infty]$, so the entire curve~\eqref{eq:farey-endpoints} occurs.

For any eventually positive unbounded nondecreasing $f$, there is a Liouville parameter in the window with $\liminf_{\eps\downarrow0}D_\lambda(\eps)/f(1/\eps)=0$. Indeed choose a continued-fraction prefix in the window, then at stage $j$, with present denominator $q$, choose the next denominator $Q$ so large that
\[
 f(q^3Q/C)\ge20jq,\qquad Q\ge q^j.
\]
Later tails preserve the first-intermediate product bound up to the same absolute constant. At $N=q+q_{\rm prev}\le2q$ and actual deficit $\eps_j=g(N)$, the dimension is at most $20q$ while $1/\eps_j\ge q^3Q/C$. The ratio is at most $1/j$. The second condition makes the parameter Liouville. The parameter depends on $f$, and this is a subsequence conclusion; its upper exponent remains $1/2$.

The exact-marginal problem does not by itself decide any dimension floor for Boolean value-one games. In fact for a fixed Boolean predicate and finite support, any strictly positive real question distribution $p$ is comparable to the rational uniform distribution $u$: if $c=|T|\min p_t$ and $C=|T|\max p_t$, then $cL_u\le L_p\le CL_u$. Constant precision rescaling preserves polynomial exponents. Thus a claim that the probability rationality alone explains a floor would not follow from Theorem~\ref{thm:farey}.
